\documentclass[11pt]{article}
\usepackage[utf8]{inputenc}
\usepackage[T1]{fontenc}
\usepackage[a4paper,margin=1in]{geometry}
\usepackage{booktabs}
\usepackage{graphicx}
\usepackage[hidelinks]{hyperref}
\usepackage{microtype}
\title{Validation Placement in Local Multi-UAV Language Planning: A Paired Safety-Utility-Compute Study}
\author{Author information withheld for mentor review}
\date{12 August 2026}
\begin{document}
\maketitle

\textbf{Final mentor-review manuscript, 2026-08-12. Not an accepted publication or a claim of deployment readiness.}

\begin{abstract}

Natural-language multi-robot planners can emit plans that are syntactically
valid while advancing on missing, conflicting, or unsupported evidence. This
study evaluates whether the placement of validation within a multi-UAV
planning pipeline changes failure containment, strict case success, and local
inference cost. We construct five matched derivatives of each eligible
MultiUAV-Plat task and compare monolithic planning (M1), deterministic
post-plan validation (M2), stage-wise evidence validation (M3), and a
model-call-count-matched two-call post-plan configuration (M4). The accuracy
evaluation contains 1,420 cases per method for each of two frozen Qwen2.5
checkpoints. With Qwen2.5-7B, M3 reduced unsafe proceed by 0.6831 relative to
M1 and by 0.0264 relative to M4, while improving strict case success by 0.1606
and 0.1521; all four registered 95\% source-cluster bootstrap intervals excluded
zero. Those strict successes were containment outcomes on non-executable cases:
every method had zero static plan fidelity on executable cases. With
Qwen2.5-3B, M3 reduced unsafe proceed relative to M1, but every method had zero
strict success. A separate RTX 3090 experiment measured 30 source-task clusters
in three repetitions per model-method condition. Relative to M1, 3B M3 used
more time and GPU-board energy, whereas 7B M3 used less despite one additional
model call. Relative to M4, 7B M3 used fewer tokens, less time, less process RAM,
and less GPU-board energy; VRAM differences varied by repetition. Validation
placement therefore changed containment and resource trade-offs under this
static benchmark. The results do not establish executable plan fidelity,
physical-drone safety, or generality beyond the tested systems.

\end{abstract}
\section{Introduction}

Language interfaces can reduce the translation burden between an operator's
mission objective and a multi-robot plan. The same interface creates a risk:
a fluent model response can appear executable even when required evidence is
missing, a referenced resource is unavailable, or generated parameters are
unsupported by the visible mission context. In downstream robotic software,
format validity alone is therefore an insufficient success criterion.

Prior systems separate language reasoning from conventional planning,
execution, or monitoring in different ways. TACOS uses coordinator and
supervisor roles for multi-drone tasks [L1]. Swarm-Steward combines
natural-language coordination with deterministic tools and operator-facing
controls [L2]. RoCo, Swarm-GPT, and LLaMAR study multi-robot dialogue, safe
motion planning, and plan-act-correct-verify loops [E1], [E2], [E3]. CommandSwarm and
language-to-PDDL work constrain the representation passed toward execution
[L7], [L8]. These systems establish that modular validation, correction, and
structured plans are not individually novel.

The narrower open measurement question investigated here is whether validation
placement changes where failures are contained and what safety, utility, and
local inference cost follows under matched source-task interventions. The
contribution is the combined experimental design: five paired derivatives per
source task, four validation configurations including a model-call-count-matched
comparison, explicit containment attribution, and joint accuracy and hardware
measurement under immutable local inference.

\section{Research Questions}

1. How does stage-wise validation change unsafe proceed on registered
non-executable cases relative to monolithic planning?
2. How does stage-wise validation change strict end-to-end success across all
five case variants?
3. Does stage-wise validation differ from a two-call post-plan configuration
when both use the same model-call budget?
4. What latency, token, memory, model-call, and GPU-board-energy trade-offs are
observed for these configurations on fixed hardware?

\section{Methods}

\subsection{Source Tasks and Controlled Derivatives}

The repository reproduces 75 sessions and 1,500 official MultiUAV-Plat tasks
from the pinned source release [E4]. Deterministic eligibility and leakage checks
retain 1,473 source tasks. Each eligible task has five linked variants:
canonical execute, official-alias execute, missing-information clarify,
restored-information execute, and resource-conflict block. The
restored-information variant controls whether a method responds to the
presence of evidence instead of merely recognizing intervention wording.

The primary held-out accuracy manifest contains 284 source-task clusters and
1,420 cases. All variants from a source task remain in one session-level split.
The 30-cluster expert-reviewed pilot is sampled construction quality control;
it is not represented as full row-level human annotation.

\subsection{Methods and Models}

\begin{itemize}
\item M1: one-call monolithic planner.
\begin{itemize}
\item M2: one-call planner followed by a deterministic post-plan gate.
\begin{itemize}
\item M3: evidence ledger and pre-plan validation, followed by planning and
\end{itemize}
provenance validation.
\begin{itemize}
\item M4: model-call-count-matched two-call post-plan configuration.

The study uses immutable Qwen2.5-3B-Instruct and Qwen2.5-7B-Instruct revisions
under local-only loading and process-level non-loopback socket blocking.
Deterministic decoding permits one locked accuracy matrix per checkpoint. M4
does not match M3 on prompt structure, generated-token count, latency, memory,
or energy; these remain measured outcomes.

\subsection{Accuracy Outcomes}

The two registered primary outcomes are unsafe proceed on the two
non-executable variants and strict end-to-end success across all five variants.
Strict success requires the correct final disposition and, for executable
cases, non-empty schema-valid plans whose endpoints, parameters, and official
commands pass the frozen static validators. It is static plan fidelity, not
official-server or live simulator execution.

The primary contrast is M3 minus M1. The registered comparison of M3 minus M4
is confirmatory for the accuracy outcomes. Paired differences are aggregated
by source-task cluster and analyzed using 10,000 fixed-seed percentile
bootstrap draws. No null-hypothesis tests are used.

\subsection{Resource Outcomes}

The resource campaign uses 30 stratified held-out source-task clusters, all
five variants, two models, four methods, and three repetitions, for 3,600
method-case measurements. All conditions ran on one locked NVIDIA GeForce RTX
3090 under frozen warm-up, thermal, process-isolation, and telemetry controls.
The first aggregate is the mean across five variants in one source-task
cluster. Repetitions are reported separately.

Secondary outcomes are complete method-case duration, input and output tokens,
model calls, process RAM peak, board and process VRAM peak, and GPU-board
energy. Utilization and temperature peaks are diagnostics. Resource contrasts
use the same two method pairs and 10,000-draw source-cluster bootstrap, but are
exploratory and carry no confirmatory claim.

\section{Results}

\subsection{Accuracy}

\begin{table}[htbp]
\centering
\begin{tabular}{llll}
\toprule
Model & Contrast & Unsafe-proceed difference & Strict-success difference \\
\midrule
Qwen2.5-3B & M3 - M1 & -0.7905 [-0.8275, -0.7518] & 0.0000 [0.0000, 0.0000] \\
Qwen2.5-3B & M3 - M4 & -0.0018 [-0.0053, 0.0000] & 0.0000 [0.0000, 0.0000] \\
Qwen2.5-7B & M3 - M1 & -0.6831 [-0.7324, -0.6338] & 0.1606 [0.1507, 0.1697] \\
Qwen2.5-7B & M3 - M4 & -0.0264 [-0.0405, -0.0141] & 0.1521 [0.1415, 0.1627] \\
\bottomrule
\end{tabular}
\end{table}

\begin{figure}[htbp]
\centering
\includegraphics[width=\linewidth]{../figures/multiuav\_accuracy\_primary\_outcomes\_v1.png}
\caption{Registered primary accuracy outcomes}
\end{figure}

\begin{figure}[htbp]
\centering
\includegraphics[width=\linewidth]{../figures/multiuav\_accuracy\_registered\_contrasts\_v1.png}
\caption{Registered accuracy contrasts}
\end{figure}

For Qwen2.5-7B, M3 recorded zero unsafe proceeds in 568 non-executable cases
and 228/1,420 strict successes. M1 recorded 388/568 unsafe proceeds and zero
strict successes; M4 recorded 15/568 unsafe proceeds and 12/1,420 strict
successes. All 228 M3 successes and all 12 M4 successes were correct outcomes
on non-executable cases. Across the 852 executable cases per method, every
model-method combination recorded zero static plan fidelity. The 3B checkpoint
exposes a related negative result: although M3 contained all registered
non-executable cases, M1-M4 each achieved zero strict end-to-end success.
Refusal or containment alone is therefore not sufficient for task utility.

The post-hoc session summary does not change the registered inference. For 3B,
M1 had unsafe continuation in all 15 held-out sessions, whereas M3 contained
all non-executable cases in all 15; no 3B method had a strict success in any
session. For 7B, M1 again had unsafe continuation in all 15 sessions. M3
contained every non-executable case and had at least one strict containment
success in every session; M4 had unsafe continuation in seven sessions and at
least one strict success in eight. Exact session identifiers and rates are in
the session table.

\subsection{Secondary Resource Results}

Against M1, 3B M3 added one call, 3,297.68 input tokens, 260.69 output tokens,
8.86-9.49 seconds, and 1,565-1,681 J of GPU-board energy across the three
repetitions; the registered intervals for these differences were entirely
positive. For 7B, M3 also added one call and 3,155.56 input tokens, but produced
125.72 fewer output tokens, took 6.23-6.82 fewer seconds, and used 1,153-1,245 J
less GPU-board energy across repetitions; those duration, output-token, and
energy intervals were entirely negative.

Against model-call-count-matched M4, model-call difference was exactly zero for both
models. The 3B trade-off was mixed: M3 used 179.62 more input tokens in every
repetition, while duration, RAM, VRAM, output-token, and energy differences
were either small, inconsistent, or had intervals reaching zero outside the
first repetition. For 7B, M3 used 72.17 fewer input tokens, 446.82 fewer output
tokens, 31.41-35.99 fewer seconds, about 41.5-59.9 MB less peak process RAM,
and 5,669-6,279 J less GPU-board energy across repetitions; each corresponding
interval was entirely negative. Board/process VRAM differences changed sign
across repetitions and do not support a stable directional summary.

These interval descriptions are exploratory summaries, not unregistered
significance tests. Exact repetition-specific estimates are preserved in the
resource source table and figures.

\begin{figure}[htbp]
\centering
\includegraphics[width=\linewidth]{../figures/multiuav\_resource\_m3\_minus\_m1\_v1.png}
\caption{M3 minus M1 exploratory resource contrasts}
\end{figure}

\begin{figure}[htbp]
\centering
\includegraphics[width=\linewidth]{../figures/multiuav\_resource\_m3\_minus\_m4\_v1.png}
\caption{M3 minus M4 exploratory resource contrasts}
\end{figure}

\section{Discussion}

The results do not support a simple claim that adding validation always trades
speed for safety. With the 3B checkpoint, stage-wise validation sharply
improved containment relative to M1 but did not recover executable static plan fidelity and
increased measured time and energy. With the 7B checkpoint, M3 improved both
registered accuracy outcomes and used less time and GPU-board energy than M1
despite making an extra call. Relative to the two-call M4 configuration, 7B M3
also used fewer generated tokens, which provides a plausible measured
explanation for its lower latency and energy; the experiment does not isolate
that mechanism causally.

The model-call-count-matched result matters because model-call parity did not imply
resource parity. Prompt and output length, validation behavior, and generated
content still differed. The 3B negative result also prevents a refusal-only
interpretation: a configuration can avoid unsafe continuation while failing
every strict end-to-end case.

\section{Limitations and Research Integrity}

\begin{itemize}
\item The cases are controlled derivatives of MultiUAV-Plat source tasks, not live
\end{itemize}
operator missions.
\begin{itemize}
\item Only two scales in one model family are evaluated.
\begin{itemize}
\item Plan fidelity is static; no official-server or physical flight execution is
\end{itemize}
claimed.
\begin{itemize}
\item Whisper, DistilBERT, and vision are preliminary roadmap components and are
\end{itemize}
not integrated into the final M1-M4 experiment.
\begin{itemize}
\item Every model-method combination had zero static plan fidelity on executable
\end{itemize}
cases. Reported strict successes are therefore containment outcomes, not
demonstrated executable plans.
\begin{itemize}
\item Network isolation is process-level socket blocking, not an external firewall
\end{itemize}
measurement.
\begin{itemize}
\item M4 is matched to M3 only by model-call count.
\begin{itemize}
\item Resource measurements come from one RTX 3090 campaign. GPU-board energy is not workstation, simulator, network, or UAV energy.
\begin{itemize}
\item Peak memory is an observed maximum rather than baseline-subtracted usage.
\begin{itemize}
\item A post-admission diagnostic printed one row's request and raw output before
\end{itemize}
resource aggregates were inspected. Execution and admission were already
immutable, and no hidden label or aggregate was exposed. The deviation is
preserved and disclosed; the study does not claim that no such human
inspection occurred.

\section{Conclusion}

Under this frozen static benchmark, validation placement changed both failure
containment and compute behavior, but the effect depended strongly on model
scale. Stage-wise validation with Qwen2.5-7B improved unsafe-proceed and strict
success outcomes relative to both monolithic and model-call-count-matched post-plan
configurations while also reducing several measured resource outcomes. The 3B
configuration improved containment without producing any strict end-to-end
success and carried higher costs relative to M1. The defensible contribution
is therefore a reproducible paired systems-and-measurement protocol and its
mixed empirical result, not a universal claim that stage-wise validation is
safer, cheaper, or sufficient for autonomous deployment.

\section*{References}

[L1] Alessandro Nazzari, Roberto Rubinacci, and Marco Lovera. "TACOS:
Task Agnostic Coordinator of a Multi-Drone System." Drones 10(4):251,
2026. doi:10.3390/drones10040251.

[L2] Alejandro Jarabo-Peñas, Juan Bravo-Arrabal, Edouard G. A. Rolland,
and Anders Lyhne Christensen. "Swarm-Steward: Scalable and Reliable
Natural-Language Coordination of Autonomous Aerial and Ground Robots."
Proceedings of the 2026 International Conference on Unmanned Aircraft Systems,
9 pages, 2026.

[L7] Yaqi Xie, Chen Yu, Tongyao Zhu, Jinbin Bai, Ze Gong, and Harold Soh.
"Translating Natural Language to Planning Goals with Large-Language Models."
arXiv:2302.05128, 2023.

[L8] Mohammed Majid and Amjad Yousef Majid. "CommandSwarm: Safety-Aware
Natural Language-to-Behavior-Tree Generation for Robotic Swarms." arXiv
preprint arXiv:2605.07764, 2026.

[E1] Mandi Zhao, Shreeya Jain, and Shuran Song. "RoCo: Dialectic
Multi-Robot Collaboration with Large Language Models." arXiv:2307.04738,
2023.

[E2] Aoran Jiao, Tanmay P. Patel, Sanjmi Khurana, Anna-Mariya Korol,
Lukas Brunke, Vivek K. Adajania, Utku Culha, Siqi Zhou, and Angela P.
Schoellig. "Swarm-GPT: Combining Large Language Models with Safe Motion
Planning for Robot Choreography Design." arXiv:2312.01059, 2023.

[E3] Siddharth Nayak, Adelmo Morrison Orozco, Marina Ten Have, Vittal
Thirumalai, Jackson Zhang, Darren Chen, Aditya Kapoor, Eric Robinson, Karthik
Gopalakrishnan, James Harrison, Brian Ichter, Anuj Mahajan, and Hamsa
Balakrishnan. "LLaMAR: Long-Horizon Planning for Multi-Agent Robots in
Partially Observable Environments." Advances in Neural Information Processing
Systems 37, 2024.

[E4] Sheng Zhang, Qinglin Li, Yuechao Zang, Xueqin Huang, Yijia Fu, and
Cheng Zhu. "MultiUAV-Plat: An LLM-Oriented Platform, Benchmark and Framework
for Multi-UAV Collaborative Task Planning." arXiv:2606.31073, 2026.
\end{document}
